\chapter{Falsification, Revolutions, and Research Programmes}
\label{ch:change}

\begin{chapterguide}
This chapter moves from the logic of evidence to the history of scientific
change. Popper's falsificationism, Kuhn's account of normal and revolutionary
science, Lakatos's research programmes, and Feyerabend's methodological
pluralism are treated as complementary diagnoses of different pressures within
science. The central question is whether historical change destroys rational
progress or reveals a more complex form of it.
\end{chapterguide}

The logical empiricist image of science was often read as a picture of
hypotheses confronting observations one by one. Karl Popper objected that this
picture misunderstands the asymmetry between confirmation and criticism.
Universal statements cannot be conclusively verified by finite observations,
but a universal statement can conflict with a particular observation once the
auxiliary conditions are fixed. Popper therefore emphasized conjectures,
deductive consequences, and attempts at refutation \citep{popper1959}.

Later historians and philosophers showed that scientific practice does not fit a
simple rule of immediate abandonment. Anomalies are common. Instruments fail.
A theory may be central to a research programme and worth preserving while an
auxiliary assumption is revised. Scientific communities work within inherited
standards and problem-solutions. The important task is not to choose between
logic and history but to understand how criticism operates inside historical
practice.

\section{Popper's strongest falsificationism}

Popper's positive thesis is not that scientists can mechanically prove theories
false. A scientific theory is exposed to risk when it rules out possible
observations and when its predictions are precise enough to be vulnerable. A
theory that accommodates every possible outcome by changing its interpretation
after the fact is not thereby scientific.

This is a criterion of \emph{testability}, not a complete description of
scientific careers. It explains why a bold theory can be epistemically
valuable even before it is established: it tells us what the world would have
to be like for the theory to fail. It also explains why ad hoc adjustments are
suspicious when they merely protect a claim from evidence without generating
new constraints.

Popper's account remains useful for demarcation and scientific honesty. It is
not enough to collect confirming examples. A serious claim should expose
itself to possible failure. But two qualifications are necessary.

First, observations do not test isolated theories. A failed prediction can
reflect an instrument error, a hidden confounder, or a mistaken initial
condition. Second, a theory may rationally survive an anomaly when it is part
of a productive programme whose other consequences are successful. Immediate
falsification is too crude as a rule of practice.

\section{Kuhn and the structure of scientific change}

Thomas Kuhn's account of scientific revolutions begins with a historical
observation: much of mature science is not a continuous search for decisive
tests of foundational principles. Scientists work within a paradigm, solving
puzzles using shared exemplars, standards, instruments, and concepts
\citep{kuhn1962}. This normal science is conservative in a productive sense.
It concentrates effort and permits detailed work that would be impossible if
every assumption were reopened at every moment.

A paradigm can nevertheless generate persistent anomalies. When the anomalies
interact with practical failures, new instruments, or a rival conceptual
framework, a period of crisis may arise. A revolution changes not only an
isolated theory but the questions that count as important, the objects that
are salient, and the standards by which solutions are judged.

The strongest reading of Kuhn is not that scientists live in private worlds
with no contact across paradigms. Kuhn explicitly retained a notion of
progress in puzzle-solving and argued that later theories often extend the
range and precision of problems that can be addressed. His discussion of
incommensurability concerns difficulties of translation and comparison, not a
claim that all theories are equally good.

Kuhn's insight is that standards are partly embodied in practice. A criterion
such as simplicity or accuracy does not decide theory choice without a
background understanding of what counts as a relevant problem and a workable
representation. The danger is that this insight can slide into relativism if
the shared world disappears from the account. Anomalies are not merely
failures to fit a community's expectations; they can be persistent effects
generated by the world.

\section{Lakatos and the time-scale of criticism}

Imre Lakatos attempted to retain Popperian rationality while acknowledging the
historical endurance of research traditions. A research programme has a
\enquote{hard core} of central commitments, a \enquote{protective belt} of
auxiliary hypotheses, and heuristic rules guiding the development of the
programme \citep{lakatos1970,lakatos1978}. Scientists do not abandon a
programme after every anomaly. They compare programmes over time.

A programme is progressive when modifications produce new empirical content
and some of that content is corroborated. It is degenerating when changes
primarily accommodate known facts without leading to novel success. The
standard is comparative and historical. A theory is not condemned by a single
bad result, but a programme can lose rational support if a rival repeatedly
generates and confirms more demanding results.

Lakatos improves both Popper and Kuhn. He captures the importance of research
continuity without treating persistence as irrational, and he turns the vague
idea of paradigm success into a comparative methodological question. But the
criterion is not mechanically decisive. The distinction between a novel
prediction and a post hoc accommodation can depend on how a research history
is reconstructed. A degenerating programme can recover after a long pause.
Science also contains multiple programmes that address different targets and
cannot be ranked on a single scale.

\section{Feyerabend and the criticism of methodological monism}

Paul Feyerabend argued that no fixed methodological rule captures the history
of major scientific advances. The case of Galileo, for example, involved
conceptual, instrumental, rhetorical, and experimental strategies that would
look illegitimate if judged by a narrow rule of confirmation. Feyerabend's
provocation was directed not only at method but at the political authority that
can arise when one method is treated as a monopoly \citep{feyerabend1975}.

The charitable interpretation is pluralist rather than nihilist. Scientific
progress can require changing the standards by which a problem is approached.
Alternative theories can reveal the hidden assumptions of a dominant theory.
A period of methodological diversity may be epistemically useful because it
prevents a community from becoming blind to its own categories.

The strongest objection to Feyerabend is that \enquote{anything goes} cannot
itself be a sufficient norm. Some practices are more transparent, more
reliable, and more capable of correcting error than others. A community that
suppresses criticism or protects a claim from every possible test is not
rescued by historical pluralism. \RCR\ accepts methodological diversity but
asks what constraints the method creates and whether those constraints are
genuinely world-sensitive.

\section{Are revolutions rational?}

The phrase \enquote{scientific revolution} can suggest a discontinuity so
complete that rational comparison becomes impossible. There is a real issue
here: concepts change their meaning when their inferential role changes. The
term \enquote{mass} in Newtonian mechanics is not simply the same term as
relativistic mass. Some standards shift with the problem situation.

Yet complete discontinuity is too strong. A later theory may preserve limiting
relations, numerical results, experimental techniques, or causal interventions
from an earlier one. Old concepts may survive as approximations within a
restricted domain. Even when vocabularies change, laboratories, bodies,
detectors, and failed predictions can constrain both sides.

\begin{principlebox}
\textbf{Historical continuity principle.} Rational comparison across theories
does not require a vocabulary-free standpoint. It requires enough shared
constraints that rival representations can be linked to common phenomena,
operations, or consequences, even when their concepts are not fully
intertranslatable.
\end{principlebox}

This principle also clarifies why scientific progress is not mere accumulation.
A later theory can improve a representation by changing the dimensions along
which the phenomenon is understood. Progress may be conceptual, predictive,
causal, experimental, or organizational. It is not a single scalar quantity.

\section{The lesson for objectivity}

These traditions yield a layered picture of scientific rationality.

\begin{table}[ht]
\centering
\smalltable
\begin{tabularx}{\textwidth}{p{0.22\textwidth}YY}
\toprule
Tradition & Key achievement & Required correction\\
\midrule
Popper & Risk, criticism, testability & Tests involve auxiliary assumptions and time\\
Kuhn & Practice, paradigms, conceptual change & Revolutions remain constrained by shared phenomena\\
Lakatos & Comparative research programmes & Progress is plural and historically underdetermined\\
Feyerabend & Methodological pluralism and criticism of monopoly & Pluralism needs standards of transparency and error correction\\
\bottomrule
\end{tabularx}
\caption{Four accounts of scientific change and the constraints needed to
make them compatible.}
\label{tab:change-traditions}
\end{table}

Scientific objectivity cannot be identified with resistance to all change. A
dogmatic theory can survive anomalies. Nor can it be identified with change
itself. A community can shift because of fashion, power, or institutional
pressure. Objectivity concerns the quality of the constraints that govern
change: whether anomalies are investigated, alternatives are developed,
evidence is independently checked, and revisions increase the range and
reliability of world-directed inference.

\section*{Chapter summary}

Popper shows why risky claims and attempts at refutation matter. Kuhn shows
that scientists work within paradigms and that revolutions change problems and
standards. Lakatos explains why rational criticism often operates over the
history of research programmes rather than through instant rejection.
Feyerabend warns against methodological monopoly and reminds us that new
methods can be necessary for new problems. Together they support a fallibilist,
historical, and plural account of rationality, provided that theory change
remains constrained by common phenomena, interventions, and error correction.

\begin{discussion}
\begin{enumerate}
\item When should an anomaly be treated as evidence against a theory rather
than as a reason to revise an auxiliary assumption?
\item Can a research programme be progressive in one domain and degenerating in
another?
\item Does conceptual incommensurability prevent rational comparison, or only
make comparison more demanding?
\end{enumerate}
\end{discussion}

\begin{furtherreading}
Compare Popper (1959), Kuhn (1962), Lakatos (1970), and Feyerabend (1975)
directly. For critical discussion of historicist rationality, consult the
relevant scholarship on scientific change and the essays collected in
Lakatos and Musgrave (1970).
\end{furtherreading}

